\section[PHT3D Exercise 5: Impact of cation exchange on groundwater-bentonite interaction]{PHT3D Exercise 5: Impact of cation exchange on groundwater-bentonite interaction}

Bentonite clays are candidate materials for engineered barrier
systems for the encapsulation of hazardous wastes because of their
low hydraulic permeability, swelling potential, ability to self-heal
cracks in contact with water and their sorption potential.
%\cite{Pusch_2001} 
Compacted bentonite buffers are being investigated in
deep geological repositories for radioactive waste 
(Figure \ref{fig:nagra}) for their suitability
to encapsulate waste canisters and ability to seal access
galleries, isolating the nuclear waste from the hydrogeological
environment. Accordingly, the geochemical stability of bentonite
clays is essential for the performance and safety evaluation ofwaste
repositories. As part of these safety assessments, experiments on bentonite
durability and stability over time and under different thermal,
geochemical and mechanical conditions at different scales have
been performed over the years. Potential mineral transformations
and alterations in porewater chemistry and mineralogy under welldefined
conditions have been investigated in laboratory experiments
in order to predict solute mobility and the geochemical
stability of bentonites. During the re-saturation phase after placement
of spent nuclear fuel, the transport of solutes in bentonite
clays has been found to be controlled by advection \citep{salas_2014}.

\begin{figure}[hb]
\centering
\includegraphics[width=0.6\textwidth, angle = 0]{nagra}
%\includegraphics[height=0.3]{mt3dms_ex1_map.eps}
%\includegraphics{mt3dms_ex1_map.eps}
\caption[Schematic illustration of geological disposal concept (source: NAGRA)]{Schematic illustration of geological disposal concept (source: NAGRA)}
\label{fig:nagra}
\end{figure}

However, due to the low hydraulic conductivity of clays,
following saturation, the main transport process is in this media is by molecular diffusion. 
In order to predict solute migration under conditions expected in a
repository, the diffusive behaviour within compacted bentonites
has been investigated in laboratory experiments for different species
(e.g., water, anions, cations, colloids) 
\citep{appelo_vanloon_2010, gimmi_2011, vanloon_2007, melkior_2009, sanchez_2008} 
among others; 
under different temperatures 
\citep{kozaki_1998, suzuki_2004, sanchez_2008}
and as a function of compaction \citep{suzuki_2004}. 
%(Pusch, 2002;  

\begin{figure}[ht]
\centering
\includegraphics[width=0.6\textwidth, angle = 0]{bentonite_expt}
%\includegraphics[height=0.3]{mt3dms_ex1_map.eps}
%\includegraphics{mt3dms_ex1_map.eps}
\caption[ABM Experiment]{ABM Experiment}
\label{fig:abm}
\end{figure}

Nevertheless, the underlying phenomena controlling
diffusion within bentonites is still an active field of research,
including their mechanistic representation in numerical models.
%(Idiart and Pekala, 2016).
Uncertainties result from a lack of a well-established and widely
accepted conceptual model of the structure of compacted bentonite
at the nano-scale at different dry densities and salinities. Available
characterisation techniques for porosity distribution and nanostructure
of compacted bentonite often perturb the structure of
the water saturated bentonite and are therefore of limited applicability.
In addition, the validity of the classical Fickian diffusion
approach based on a single bulk (total) porosity has been questioned
in some cases, as it is unable to represent the diffusion of
cations, anions and neutral species simultaneously. As a result,
more complex mechanistic models have been proposed in the
literature to explain experimental data, based on different conceptualizations
of the pore space, e.g., \citep{birgersson_2009, appelo_vanloon_2010, gimmi_2011, tachi_2014}.

Solute mobility in compacted 
bentonites is further influenced by
geochemical reactions such as cation and anion exchange 
\citep{valter_2013, birgersson_2009, tachi_2014}, 
as well as the dissolution of primary minerals or the
precipitation of secondary phases \citep{arcos_2008, kaufhold_2013}, 
%(Liu and Neretnieks, 2006; 
%Idiart et al., 2013) 
which impact on the performance of bentonites as barriers 
by altering, e.g., the swelling capacity or the hydraulic and gas 
conductivities of these clays. 
In order to upscale laboratory findings and to test the
behaviour of bentonite clays under the thermo-hydraulic conditions
expected in radioactive waste repositories, several countries
are conducting large scale in-situ tests of engineered barrier systems.

In 2006 a field scale experiment was launched, the Alternative Buffer
Material (ABM) test 
%(Eng et al., 2007; 
\citep{svensson_2011}
, which extended investigations to a range of clays, including compacted
sedimentary clay materials (Friedland clay, Germany and Callovo-
Oxfordian clay, France) and various compacted bentonites. 

Three columns consisting of these different compacted clay buffer blocks
were installed in three boreholes (ABM1-3). The analysis of the first
package, ABM1 as illustrated in Figure \ref{fig:abm}, 
is the subject of this modelling exercise. 
The aim of the ABM1 test was to compare different clay materials 
and their performance over time when in contact with 
the heater, the seals, the fractured
host rock, and ambient groundwater at relevant depths as well as
the expected maximum in-situ temperatures. 
The main objective was to provide a basis for selecting an ideal 
buffer for the Swedish spent nuclear fuel repository. 
The experiment consists of a vertical borehole placed at
450 m depth \citep{svensson_2011}. 

It contains a central heater inside an iron tube surrounded by 11
different compacted clay blocks. The column of blocks was artificially
saturated with Aspo groundwater and heated to a maximum
temperature of 130 C. Over a duration of 28 months data were
recorded on the temperature and water pressure evolution.
Following the test, a comprehensive characterization of the cation
exchange population and distribution with time \citep{dohrmann_2013}
as well as mineralogical alterations \citep{kaufhold_2013} 
was carried out. All buffers started with distinctly different
exchangeable cation (EC) populations and it was expected that,
after the field experiment, the ECs would display horizontal variations
for all bentonite rings due to the radial heat distribution
induced by the central heater element. However, in contrast to this,
large scale variations in the ion exchange composition had
occurred. That is, the composition of exchangeable cations had
changed markedly in all clay materials over the whole test parcel
after 28 months, but no significant horizontal variations in the EC
populations within each ring was detectable. The most significant
change was an increase in exchangeable $Ca^{2+}$, and a decrease in
exchangeable $Na^{+}$ and Mg$^{2+}$. 
Mineralogical changes of the bentonites
were unable to explain these significant changes, and it was
postulated that interaction with the surrounding groundwater
must have caused partial equilibration of the porewater chemistry
within and possibly between the different bentonite blocks within
a short time frame \citep{dohrmann_2013}. 
%(Idiart et al., 2014; 
%Pekala et al., 2012)

In order to assess this hypothesis and to enhance the understanding
of the processes controlling the changes in porewater chemistry 
in bentonites under the influence of in-situ groundwater 
and heat conditions, a numerical model of the ABM1 
field test was developed.
For the modelling study the aim was (i) to develop a numerical model that
could provide a process-based description of the physical and
geochemical mechanisms controlling the chemical evolution of
bentonite porewaters at the field (i.e., meter) scale; 
(ii) to quantify the influence of groundwater on the chemical evolution of the
different clays with time, and (iii) to verify the model by applying it
to the available field data set collected during the ABM1 test.

The set-up of the ABM1 experiment, which is described in detail
in x 
%Eng et al. (2007) 
and \cite{svensson_2011}, 
consists of a central
steel canister with heaters, surrounded by ring shaped bentonite
blocks situated in a borehole of 300 mm diameter in crystalline
bedrock at about 450 m depth (Figure \ref{fig:abm}). 

Eleven different clay materials
(two sedimentary clays and nine bentonites) were included in
the ABM tests as compacted rings 
%(Table 1 and 
(Figure \ref{fig:abm}) 
with an outer diameter of 280 mm, and an inner diameter of 100 mm. 
Additionally, four rings were included which consisted of cloth-covered
steel-cages filled with bentonite pellets with and without additional
quartz to also test granulated material. All clay materials
were installed multiple times and the total number of material
blocks stacked on top of each other was 30 (Figure \ref{fig:abm}). 
Each block had a height of 100 mm, resulting in a total length of 3 m for the whole
ABM1 package. The gap between the ABM1 package and the crystalline
host rock was filled with sand, in which an artificial wetting
system was installed. This consisted of titanium tubes, connected to
a water tank, which supplied native 
Aspo groundwater over the
duration of the experiment. The wetting system was aimed at
simulating water bearing fractures within the host rock and supplied
ambient groundwater in addition to natural water inflow
through fractures into the deposit bore. However, the effectiveness
of the artificial wetting system to evenly supply groundwater to the
ABM1 package was open to question, as monitoring of applied
water pressure and injected water volume during the experiment
suggested leakage of injected water into the surrounding host rock
at some unknown depth within the bore \citep{svensson_2011}. 

The inflow of water through the fractures into the empty borehole was
measured at 3.6 l/hour before the test, but inflow locations were
not identified 
%(Eng et al., 2007). 
The ABM1 package was heated from the start of the experiment 
for the entire duration of the test
with exposure of the blocks to maximum temperatures of 130 C
for about one year (Figure \ref{fig:abm}). 

A heating element ran through the entire
axis of the ABM1 package with additional heaters being installed at
the top and bottom of the package to provide a relatively homogeneous
temperature distribution throughout the package length.
A chemical and mineralogical characterization of the different
clay materials was undertaken before the start of the test in order to
obtain a reference and to provide initial conditions before heating
and saturation \citep{svensson_2011}.
%Kumpulainen and Kiviranta, 2011)

After 28 months, the ABM1 experiment was terminated
and each block was sampled at five different radial distances from
the contact with the metal tube to record geochemical and
mineralogical alterations. Additionally, samples were collected at
three different vertical intervals in each block. It was found that,
apart from the EC distribution, almost all geochemical and mineralogical
alterations of the different clay materials were restricted to
a less than 5 mm wide zone from the interface between the central
steel canister and the blocks. Here, alterations included the accumulation
of Fe and organic carbon, the precipitation of carbonates
and gypsum/anhydride and the dissolution of some silicates
\citep{kaufhold_2013, wersin_2015}. 

Also, the cation exchange
capacity of the materials (CEC) showed little change compared to
the start of the test for most material blocks, with only 5 out of 30
showing a change of >5\% in CEC. 
In contrast, large scale alterations in the EC population 
had occurred in all blocks and modifications
occurred, without significant gradients over the entire horizontal
radial distance \citep{dohrmann_2013}.

\subsection[Basic numerical model setup]{Basic numerical model setup}

Based on the observed thermal and geochemical
changes during the ABM1 tests 
\citep{dohrmann_2013, kaufhold_2013, svensson_2011} 
%Idiart et al., 2014)
, 
we initially simulate a simplified subset of the processes that 
played during the experiment. We focus on the interaction between diffusive transport 
of the groundwater constituents and the cation exchange reactions that took place 
during the experiment. We exploit the radial symmetry of the experimental setup, 
which allows for considerable computational time savings compared to a fully 
three-dimensional setup.
Furthermore we simplify the problem by simulating only 3 different bentonite buffers
rather than the 11 different buffers that were included in the experiment and simulated
in \cite{wallis_2016}. A summary of the model setup is provided in Table \ref{wallis_setup}.
Note, that the circulating flow that occurred outside the bentonite rings was neglected.
Instead, the (constant) water composition of the circulating fluid was used to define 
a specified (constant) concentration boundary condition.  

\begin{itemize}

\item In ORTI, create a new model and save it in a new, separate directory.

\item Under {\color{blue}Parameters} $\rightarrow$ {\color{blue}Model} select the {\color{blue}Model} button and specify the attributes of the new model. Select the values {\color{blue}Radial} for {\color{blue}Dimension}, {\color{blue}unconfined} for {\color{blue}Type} 
and {\color{blue}MODFLOW series} for {\color{blue}Group} 

\item Under {\color{blue}Parameters} $\rightarrow$ {\color{blue}Model} select the {\color{blue}Grid} button and specify the details of the model domain and the spatial discretisation in the dialog that opens up. 
Define a (radial) transect that extends from Xmin = 0.05 to Xmax = 0.145m. 
For the vertical extent we consider Zmin = 0 to Zmax = 0.3m.  
The model domain may be discretised into grid cells of 5mm length and 1cm height. 
Note, that the model discretisation can be easily changed again later. 

\item Select {\color{blue}Parameters} and select the {\color{blue}Time} button.
Set the {\color{blue}Total Simulation Time} to 880 days and select a time {\color{blue}Step size}
of 5 days. 

\item Go to {\color{blue}Spatial Attributes} $\rightarrow$ {\color{blue}MODFLOW} $\rightarrow$ {\color{blue}LPF} $\rightarrow$ {\color{blue}lpf.8} to specify the horizontal hydraulic conductivity to 0.001 $m/day$.
Note that the entered hydraulic conductivity will not affect the results as the groundwater will not move
in this example. 

\item Go to {\color{blue}Spatial Attributes} $\rightarrow$ {\color{blue}MODFLOW} $\rightarrow$ {\color{blue}LPF} $\rightarrow$ {\color{blue}lpf.9} to specify the vertical hydraulic conductivity to 0.001 $m/day$

\item Go to {\color{blue}Spatial Attributes} $\rightarrow$ {\color{blue}MT3DMS} $\rightarrow$ {\color{blue}BTN} $\rightarrow$ {\color{blue}btn.11} and specify a uniform effective porosity of 0.45, even though the 
porosities of the various bentonites used in the ABM test varied slightly \citep{wallis_2016}.   


\item Given that there is no groundwater flow involved in this example there is in principle no need to define any hydraulic boundary conditions. However, for numerical reasons we define a specified head boundary at a single point by going to {\color{blue}Spatial Attributes} $\rightarrow$ {\color{blue}MODFLOW} $\rightarrow$ {\color{blue}BAS6} $\rightarrow$ {\color{blue}bas.3 Boundary Conditions} and setting the value to -1 (= active cell) at X = 0.00525 m and 0.006 m in vertical direction. 

\item We will also define a global initial hydraulic head for the entire domain (Background value) by going to {\color{blue}Spatial Attributes} $\rightarrow$ {\color{blue}MODFLOW} $\rightarrow$ {\color{blue}BAS6} $\rightarrow$ {\color{blue}bas.6 Initial Heads} and set the
 {\color{blue}Backg.} value to 1m. Then press {\color{blue}'OK'} to confirm the value.

\end{itemize}

Again, as there will not be any groundwater flow this input is solely made for numerical reasons.
For the same reason MODFLOW needs to be executed before the next stage of the model setup.

\begin{itemize}

\item Under {\color{blue}Parameters} $\rightarrow$ {\color{blue}Flow} use the {\color{blue}Write Button} to write all the input files required to run MODFLOW

\item Under {\color{blue}Parameters} $\rightarrow$ {\color{blue}Flow} use the 
{\color{blue}Red Arrow} button to start the MODFLOW simulation   

\end{itemize}

\begin{table}[t]
\caption{Model geometry and summary of key input parameters.}
\begin{center}
\begin{tabular}{lc} 
\toprule
Simulation type                     & steady state \\
Total simulation period\ $ (days)      $         &  880            \\
Time step length\ $ (days)      $         &  10            \\
Model \  length\ $ (m)      $         &  0.095            \\
Model \  thickness\ $(m)$                        &  0.3            \\
Grid \ spacing\ $\Delta x $ $(m)$                &    0.005            \\
Grid \ spacing\ $\Delta z $ $(m)$                &    0.01            \\
Porosity\     $            $         &  0.5           \\
Horizontal hydraulic conductivity\ $(m/day)$               &  0.001        \\ 
Vertical   hydraulic conductivity\ $(m/day)$               &  0.001         \\ 
Longitudinal dispersivity\ $(m)                 $         &  0.01          \\ 
Transverse  vertical dispersivity\ $(m)         $         &  0.01          \\ 
Diffusion coefficient\ $(m)         $         &  $5 \times 10^{-10}$          \\ 
\bottomrule
\end{tabular}
\label{wallis_setup}
\end{center}
\end{table}
    
\subsection[Diffusion model setup]{Diffusion model setup}

In the absence of groundwater flow inside the model domain molecular 
diffusion acted as the only transport process while neither advection 
nor hydrodynamic dispersion occurred. \cite{wallis_2016} found that the
temporal and spatial variations of the temperature-field 
had a significant impact on the diffusion rates.
To consider this effect they incorporated the temperature-dependency
of the diffusion coefficient into their model framework and obtained, 
as a result, a significantly improved match with
observed concentration data. For simplicity we neglect this dependency in this 
modelling exercise while generally it would need to be considered.     
 
\begin{itemize}

\item Go to {\color{blue}Parameters} $\rightarrow$ {\color{blue}Transport} 
and select the {\color{blue}Parameter} button


\item Select {\color{blue}Spatial Attributes} $\rightarrow$ 
{\color{blue}MT3DMS} $\rightarrow$ {\color{blue}DSP}, 
enter a diffusion coefficient of 4.32 $\times$ $10^{-5}$ $m^2$ $days^{-1}$ 
and then don't forget to click {\color{blue}OK} 

\end{itemize}
	

\begin{figure}[t]
\centering
\includegraphics[width=1.0\textwidth, angle = 0]{bento_zones}
%\includegraphics[height=0.3]{mt3dms_ex1_map.eps}
%\includegraphics{mt3dms_ex1_map.eps}
\caption[Position of Initial Chemistry zones]{Position of Initial Chemistry zones}
\label{fig:bento_zones}
\end{figure}

\begin{table}[b]
\caption{Solution compositions to define the initial water compositions and of the fluid circulating outside the bentonite rings.}
\begin{center}
\begin{tabular}{l c c c c}
\hline
 Aqueous      & $C_{Backgr}$            & $C_{init}$               & $C_{init}$                & $C_{init}$  \\
 Component    & $(mol\ l^{-1}_w)$       & $(mol\ l^{-1}_w)$        & $(mol\ l^{-1}_w)$         & $(mol\ l^{-1}_w)$   \\ \hline
 C(4)         & 8.600 $\times$ $10^{-4}$  & 8.890 $\times$ $10^{-4}$ & 8.890 $\times$ $10^{-4}$  & 8.890 $\times$ $10^{-4}$\\
 Ca           & 6.479 $\times$ $10^{-2}$  & 1.017 $\times$ $10^{-2}$ & 2.373 $\times$ $10^{-2}$  & 4.525 $\times$ $10^{-3}$\\
 Cl           & 2.332 $\times$ $10^{-1}$  & 2.332 $\times$ $10^{-1}$ & 2.332 $\times$ $10^{-1}$  & 2.332 $\times$ $10^{-1}$ \\
 K            & 3.220 $\times$ $10^{-4}$  & 2.146 $\times$ $10^{-3}$ & 2.213 $\times$ $10^{-3}$  & 1.339 $\times$ $10^{-3}$\\
 Mg           & 2.704 $\times$ $10^{-3}$  & 1.653 $\times$ $10^{-3}$ & 2.232 $\times$ $10^{-3}$  & 2.423 $\times$ $10^{-3}$\\
 Na           & 1.090 $\times$ $10^{-1}$  & 2.185 $\times$ $10^{-3}$ & 1.500 $\times$ $10^{-1}$  & 2.291 $\times$ $10^{-3}$ \\ 
 S(6)         & 5.100 $\times$ $10^{-3}$  & 5.100 $\times$ $10^{-3}$ & 5.100 $\times$ $10^{-3}$  & 5.100 $\times$ $10^{-3}$\\ 
 pH           &   7.330	                  &  7.479                   & 7.390                     & 7.499 \\
 pe           &   4                       &  4                       &   4                       & 4     \\ \hline
\end{tabular}
\label{ex_wallis_aqu_conc}
\end{center}
\end{table}
	

\begin{table}
\caption{Initial exchanger compositions associated with the 3 different bentonite rings.}
\begin{center}
\begin{tabular}{l c c c}
\hline
 Exchanger   & $C_{init}$ Assembl 1     & $C_{init}$  Assembl 2     & $C_{init}$  Assembl 3 \\
 Species     & $(mol\ l^{-1}_w)$        & $(mol\ l^{-1}_w)$         & $(mol\ l^{-1}_w)$   \\ \hline
 CaX2        & 2.290 $\times$ $10^{-1}$ & 6.281 $\times$ $10^{-1}$  & 3.244 $\times$ $10^{-1}$\\
 KX          & 3.768 $\times$ $10^{-2}$ & 5.444 $\times$ $10^{-2}$  & 5.458 $\times$ $10^{-2}$\\
 MgX2        & 7.636 $\times$ $10^{-2}$ & 3.690 $\times$ $10^{-1}$  & 1.081 $\times$ $10^{-1}$\\
 NaX         & 1.315  & 7.370 $\times$ $10^{-1}$  & 1.865 \\
\end{tabular}
\label{ex_wallis_ex_conc}
\end{center}
\end{table}


\subsection[Reactive transport]{Reactive transport}
	
\begin{itemize}

\item Copy the database file {\color{red} \textbf{pht3d\_datab.dat}} 
that is provided for this exercise 
into the folder that will contain all files for this exercise.

\item Go to {\color{blue}Parameters} $\rightarrow$ {\color{blue}Chemistry} and select the
{\color{blue}Import database} button. ORTI will then read and interpret 
the PHT3D database file.

\item In this next step the initial concentrations of the aqueous components need to be defined. 
The initial concentrations vary between the 3 different bentonite rings and therefore three different
water compositions need to be entered as {\color{red} \textbf{Solutions 1-3}}. 
Furthermore the composition of the water that is circulating 
outside the bentonite rings is entered as the {\color{red} \textbf{Backgr}} solution. 
The water compositions are listed in Table \ref{ex_wallis_aqu_conc}.
Note, that the {\color{blue}Solution} and other Tabs allow you to copy/paste 
concentrations values from external, e.g., Excel tables. 
All 4 solutions are entered under the {\color{blue} \textbf{Solutions}} Tab.
In that Tab first activate all relevant species by ticking the appropriate boxes for 
$pH$, $pe$,  $C(4)$, $Ca$, $Cl$, $K$, $Mg$, $Na$ and $S(6)$

\item In addition the initial concentrations of the four exchanger-species, i.e., $CaX2$, 
$KX$, $MgX2$ and $NaX$ need to be entered for {\color{red} \textbf{Assembl1}}, 
{\color{red} \textbf{Assembl2}} and {\color{red} \textbf{Assembl3}}, respectively,
after activating the species. 
The initial exchanger compositions are listed in Table \ref{ex_wallis_ex_conc}.
Note, that the initial exchanger composition and the initial aqueous composition 
provided in Table \ref{ex_wallis_ex_conc} are already in geochemical equilibrium.   
That implies that in the absence of transport processes no reaction would proceed.      

\item In the next step we need to define where the aqueous solutions 
and the exchanger compositions will be employed. 

\item First, {\color{red} \textbf{Backgr}} will be used to define a specified 
concentration boundary condition at the outer fringe of the bentonite ring.
This reflects the fact that the circulating water outside the bentonite rings 
maintains a constant water composition. This is implemented by 
defining the boundary type under {\color{blue}Spatial Attributes} $\rightarrow$ {\color{blue}MT3DMS} $\rightarrow$ {\color{blue}BTN} $\rightarrow$ {\color{blue}btn.12} and specify a boundary 
type as -1 (= specified concentration boundary) along a vertical line near the outer fringe (x = 0.1425).
Define a name for the boundary (such as bc\_circulation\_wq) in the dialogue.
Secondly, we need to make sure that the {\color{red} \textbf{Backgr}} water composition
is used at this boundary. To do this we normally need to define the initial water chemistry 
at the same location (here a line) as the boundary condition type. 
However, given that the {\color{red} \textbf{Backgr}} water composition is used 
as the default water composition at all location where no other initial chemistry is defined,
there is in this case no need to explicitly define the water composition for this boundary.   

\item Finally, the different initial water and cation exchanger compositions are implemented
through defining three different {\color{blue}Initial Chemistry} zones, with each 
of the three zones corresponding to one of the 0.1m thick bentonite rings.  
Define the 3 zones by going to 
{\color{blue}Spatial Attributes} $\rightarrow$ 
{\color{blue}PHT3D} $\rightarrow$ 
{\color{blue}PH} $\rightarrow$ {\color{blue}ph.3} and define one after another 
of the 3 {\color{blue}Initial Chemistry} zones (see Figure \ref{fig:bento_zones}). 
For the uppermost zone (first bentonite ring) allocate a {\color{blue}Zone Value} of {\color{red} \textbf{1010}}, 
allocate {\color{red} \textbf{2020}} for the central zone (second bentonite ring) and allocate 
{\color{red} \textbf{3030}} for the zone in which the bottom bentonite ring is located.
Setting the {\color{blue}Zone Value} to 3030 means that in this zone the chemistry is defined
through the aqueous {\color{red} \textbf{Solution 3}}, {\color{red} \textbf{Mineral assemblage 0}},
exchanger composition {\color{red} \textbf{Assembl 3}} and Surface {\color{red} \textbf{Site 0}}.
No mineral or surface complexation reactions are considered in the reaction network,
which is implemented by not activating any mineral phases or surface sites in the respective Tabs.  
 
\item Now run PHT3D.

\item Inspect the results by plotting concentration profiles of $NaX$ and $CaX2$ for 
specific simulation times. 

\item What would happen if you plotted the sum of $NaX$, $CaX2$, $KX$ and $MgX2$ ?

\item Plot aqueous and exchanger species concentrations at selected distances from the heater element and interpret the breakthrough curves as wel as the differences between the 3 bentonites.   

\end{itemize}	







